\section{Non-Functional Requirements}

In this section, there will be defined the Non-Functional Requirements that the system must achieve and support. The metrics
where used from the following decription of Airbnb's Workload Model. For its simulation, values will vary, as it is just a
``shorter'' version of the application, which will take into account its basic funcionalities.

\begin{longtable}{|p{0.12\textwidth}|p{0.70\textwidth}|}
    \hline
    \textbf{ID} & 
    \textbf{Requirement} \\
    \hline
    \endfirsthead%

    \hline
    \textbf{ID} & 
    \textbf{Requirement} \\
    \hline
    \endhead%

    NFR-01 &
    The system shall support more than 220 countries and regions, with approximately 23,000 concurrent active users under 
    normal operating conditions and approximately 60,000 concurrent active users during peak conditions. \\
    \hline

    NFR-02 &
    The system shall support an average workload of approximately 934,000 requests per hour under normal operating 
    conditions. \\
    \hline

    NFR-03 &
    The system shall support a peak workload of approximately 9.3 million requests per hour during peak operating 
    periods. \\
    \hline

    NFR-04 &
    The system shall sustain peak workloads for periods of approximately 5 minutes without degradation below the defined 
    performance and availability targets. \\
    \hline

    NFR-05 &
    The system shall process at least 95\% of user-facing requests within approximately 1.5 seconds at the 95th percentile 
    under normal operating conditions. \\
    \hline

    NFR-06 &
    The system shall maintain an availability of at least 99.9\%, corresponding to no more than approximately 43.2 minutes 
    of unavailability per calendar month. \\
    \hline

    NFR-07 &
    The system shall support an initial data volume of approximately 1 GB and accommodate an estimated annual data growth 
    of approximately 8--10\%. \\
    \hline

    NFR-08 &
    The system shall accommodate annual data growth of approximately 8--10\% of the preceding year's stored data volume 
    without requiring architectural changes to support the projected workload. \\
    \hline

    NFR-09 &
    The system shall serve users across more than 220 countries and regions and shall distribute its infrastructure and 
    traffic-handling mechanisms geographically to minimize latency for users in North America, EMEA, Latin America, and 
    Asia Pacific. \\
    \hline

    \caption{Non-Functional Requirements}
    \label{tab:nfr}

\end{longtable}

The Airbnb values in the previous table represent the scale of the real platform, so they would be unnecessarily large 
for a small academic simulation. For this project, the workload will be scaled down to 100 registered users, with 
approximately 10 concurrent active users under normal conditions and 25 during peak conditions. The simulation will 
generate approximately 500 requests per hour during normal operation and up to 1,200 requests per hour during peak 
periods, with peak periods modeled as lasting approximately 5 minutes. The response-time target will be 95th percentile 
below 2 seconds, while the system will target 99\% availability, which is more appropriate for a small local or academic 
deployment than Airbnb's production-scale availability target.

For data and geographic scope, the simulation will initially contain approximately 100 MB of application data, with an 
expected growth of approximately 20\% per year as new users, listings, reservations, and related records are added. The 
system will be designed to serve a single geographic region, rather than reproducing Airbnb's worldwide distribution, and 
therefore no geographically distributed infrastructure will be required. These values provide a manageable workload for the 
project while preserving the same types of constraints represented by the original Airbnb workload model: user concurrency, 
request volume, peak behavior, response time, availability, data growth, and geographic scope.